\section{Experiments}
\paragraph{Design.}
We fixed seven new UCI sources, their transformations and the evaluation
rules before fitting. Four tasks are binary classification and three are
regression (Table~\ref{tab:panel}). Each fit used a 100-tree random forest
of maximum depth four. Approximately 48\% of rows fit the model, 16\%
determined the deletion order, 16\% selected the support and 20\% measured
performance. Source-specific splits respected declared groups, identical
input vectors and, for Seoul and Beijing, chronology. Those two sources
used five forest seeds on the same four time periods. The other sources
used five seeded partitions. In both prospective experiments, forecasts were
sealed before test labels were released to the scorers. The UCI identifiers in
table order are 105, 563, 225, 573, 560, 501 and 440. The archive contains the source records,
row assignments, sampling caps, seeds and hashes.

On the acceptance partition, every candidate faced an intact-data check.
Macro AUC and accuracy could each fall by at most .01 relative to the full
forest. For regression, $R^2$ could fall by .02 and MAE could increase by
$.02$ times the fitting-sample response standard deviation.
The protocol called the latter a development-scale check; the code used
the fitting rows for all methods. For cross-source reporting we scale
outer MAE by the development standard deviation. Table~\ref{tab:panel}
reports MAE in each source's response units.

To separate ranking from acceptance, we evaluated the impurity path
(MDI), greedy paths using the maximum, mean or intact-size loss, and a
static single-deletion ranking. We also paired each of the four alternative
paths with MDI. Each candidate family faced four acceptance rules:
intact predictive metrics alone, or those metrics plus the intact-size,
mean-size or worst-size loss quota. The full forest completed the 37 arms.
We first average paired differences or log ratios over seeds within a
source, then weight sources equally. Brackets give descriptive 95\%
percentile intervals obtained by resampling whole sources. Repeated seeds
share observations and the seven applications form a convenience panel.

\begin{table}[t]
\centering\small
\caption{Mean columns retained and test performance over five seeds. Both
selection methods use worst-size acceptance. Scores are AUC / accuracy
in the first four rows and $R^2$ / MAE in the last three. MinPI includes
the MDI path.}\label{tab:panel}
\begin{tabular}{lrrrrr}
\toprule
Source & Inputs & MDI & MinPI & Full forest & MinPI\\
\midrule
\tableinput{minpi-predictive-rows.tex}
\bottomrule
\end{tabular}
\end{table}

\paragraph{Compression and prediction.}
The selected forests retained a geometric mean of 52.3\% of the original
columns. Our prespecified target was a further 10\% reduction against MDI
under the same acceptance rule. The observed reduction was 6.5\% (width
ratio $.935$ [$ .861,1.000$]). MinPI returned a narrower set in seven runs,
four in Beijing and three on liver patients. Ten supports exceeded the held
quota against nine for MDI. MinPI's fit-inclusive time was
$1.247$ times MDI's [$1.144,1.361$]. Adding a mean-size path to MDI was
as competitive: MinPI's relative width was $1.009$ [$ .990,1.030$].
The average changes from the full forest were $-.0054$ AUC and $-.0009$
accuracy. Regression lost $.0358$ in $R^2$ and added $.0241$ in scaled MAE.
The five Seoul fits all exceeded the held-out loss quota. Liver patients
lost $.014$ AUC while gaining $.007$ accuracy.

\paragraph{What does the minimum change?}
Holding both candidate paths fixed, switching acceptance from the mean
to the worst size reduced test quota violations from 14 to 10 of 35,
at a mean cost of $.229$ retained columns. The difference in violation
rate was $-.114$ [$-.229,-.029$]. Only three sources changed, all in the
same direction (two-sided exact sign test $p=.25$). The source-average
positive part of the worst normalized excess fell by $.0041$.
An intermediate size attained the maximum in 15 runs. For House voting,
seed 42, the intact-size excess was $.01637$ but size 11 of 16 reached
$.02512$. This was the only run where the held-out minimum-rule support
passed the intact-size quota and failed at an interior size. These
comparisons measure different consequences of the rule: changing
acceptance can change the selected support even when its eventual
violation is visible with all inputs available.

\paragraph{Computing the votes.}
In a separately frozen timing experiment, every solver started with the
same fitted forest and performed the same one-path search on the voting
sample. We paired permutations with their reversals, shared each bank
across candidates, and compared against a native sparse implementation.
At 128 permutations exact integration took $.770$ times the selector time
[$.690,.852$], including coefficient construction, and $.974$ times the
fit-inclusive time. It matched 34 of 35 sampled supports and 27 complete
paths. With eight permutations a direct implementation that avoided the
Gram matrix was faster on five of seven seed-42 runs and recovered all
seven final supports. Exact search and each of the three sampling budgets
exceeded the test quota on the same 13 of 35 runs. Exact integration removes
calibration integration error; this experiment gives no evidence of
better test predictions from removing that error.

To isolate early rejection, we also timed native kernels on five fits for
each of 25 sources (the seven above and 18 previously examined). Screening
evaluated 41.0\% of the size dot products. With allocation-free coefficient
updates, the full ranking took $.961$ times the time of a straight native
loop over all sizes [$ .923,.991$]. The corresponding ratios were $.867$
above 50 inputs and $.999$ on the seven fresh sources, all of which had
at most 16 inputs. Thus few evaluated size votes need not translate to
large wall-clock savings. The implementations returned identical orders.

\paragraph{Search and earlier benchmarks.}
After scoring we exhaustively minimized the voting-sample worst excess
at each width on the seven small sources. Greedy search missed the optimum
by more than $10^{-8}$ of the null loss at 13 widths across 12 fits. The
largest gap was $.00750$. Substituting one optimal support per width for
the greedy path, then applying exactly the same independent acceptance
rules alongside MDI, reproduced all 35 selected supports. This diagnostic
does not optimize the joint acceptance checks over all supports.
The corresponding unrestricted empirical quota problem is already
NP-hard for a fixed budget (Appendix~\ref{app:hardness}).

On the six previously studied preprocessed datasets (30 fits, up to 100
columns), the same forest procedure gave a width ratio of $.932$ and
a fit-inclusive time ratio of $1.491$ against MDI. On Madelon it retained
40.0 columns on average against 53.6, with AUC $.724$ for both. On these
same test rows, the archived gain-ranked XGBoost retrained on a selected
subset retained only 20.4 columns and reached $.830$ AUC. This comparator
can use the full development set to refit its trees. Frozen-forest compression
would be a poor choice on Madelon if retraining is available. An
earlier boosted-tree version used margin distortion and split gain. Its
13-source width ratio was $.894$ against gain but $.988$ against static
ablation. Its objective and fitted models differ from the supervised
forest comparison; we keep its results separate in the release.

\section{Scope}
The strongest observed benefit of using the minimum was fewer failed
held-out checks at a small column cost. The alternative mean-size
ranking matched its compression. Calibration failed under the Seoul
time shift, and the analytic speed advantage depended on the sampling
precision and how much training cost was included. The experiments
treat columns as equally costly and generate masks by size. They measure
neither field acquisition costs nor the natural frequency of joint sensor
failures. Deeper forests can generate many overlapping path sets, and
the Gram construction then limits the practical scope of the method.
